% 学習結果説明（材料別の温度シフト GPR。05 フォルダのデータ。用語集と図の見方）
% 05_材料ごとの機械学習データと予測検証データ フォルダで「lualatex 学習結果説明_材料別温度シフトGPR.tex」を2回実行すると PDF を作り直せる（TeX Live の LuaLaTeX を使用）。
% 図は、各材料フォルダの <材料フォルダ>_学習結果/ の PNG を相対パスで読み込む。
\documentclass[paper=a4paper,fontsize=10.5pt,line_length=44zw,number_of_lines=38]{jlreq}
\usepackage[haranoaji]{luatexja-preset}
\usepackage{amsmath}
\usepackage{graphicx}
\usepackage{booktabs}
\usepackage{tabularx}
\usepackage{enumitem}
\usepackage{float}
\usepackage{xcolor}
\usepackage[hidelinks,bookmarksnumbered]{hyperref}

\newcommand{\file}[1]{\texttt{\detokenize{#1}}}
\newcommand{\gd}{\dot{\gamma}}
\newcommand{\Tff}{T_{\mathrm{ff}}}
\newcommand{\Tset}{T_{\mathrm{set}}}
\newcommand{\Eh}{E_{h}}
\newcommand{\Ev}{E_{v}}
\newcommand{\PIII}{P3\_3}
\newcommand{\PII}{P3\_2}
\newenvironment{termlist}{\begin{itemize}[label=・,leftmargin=1\zw,labelwidth=1\zw,labelsep=0pt,itemsep=0.4ex,topsep=0.6ex,parsep=0pt]}{\end{itemize}}
\newcommand{\term}[1]{\item \textbf{#1}：}
\newenvironment{points}{\begin{itemize}[label=$\triangleright$,leftmargin=1.5\zw,itemsep=0.3ex,topsep=0.5ex,parsep=0pt]}{\end{itemize}}
\newcommand{\pt}[1]{\item \textbf{#1}：}
\renewcommand{\arraystretch}{1.25}

\title{材料別の温度シフト GPR の学習結果の説明\\[0.6ex]{\large 05 のデータ（5data で学習、3data で検証）：用語集とグラフの見方（機械学習・統計を初めて学ぶ人向け）}}
\author{}
\date{@@DATE@@}

\begin{document}
\maketitle
{\small\tableofcontents}

% =====================================================================
\section{この文書について}

この文書は、\file{05_材料ごとの機械学習データと予測検証データ} のデータだけを使って作った、\textbf{温度のシフトファクターを表せるガウス過程回帰（GPR）の樹脂粘度モデル}の結果（図と表）を、機械学習や統計を初めて学ぶ人でも読めるように説明したものです。これまでの検討（本体の手順1〜5b）とは\textbf{別のプロジェクト}で、説明変数に圧力 \PIII\ と流動中樹脂温度 $\Tff$ を使う点が違います（ユーザーの指定、10 月 14 日）。

\begin{points}
\pt{対象の結果}材料ごとに、学習データ・検証データのフォルダと同じ場所の \file{<材料フォルダ>_学習結果/}（例：\file{01_バージンPP再々実験/01_バージンPP再々実験_学習結果/}）。中に、データの確認（\file{データ確認/}）と、入力の取り方ごとのフォルダ（\file{ショットごとの入力/}・\file{条件平均の入力/}）があります。図 @@N_FIG@@ 枚、CSV @@N_CSV@@ 個（3 材料・2 通りの合計）。図は、すべて同じフォルダの \file{highres/} にも 600 dpi で保存されます。4章の図のファイルの場所は、\file{<材料フォルダ>_学習結果/} からの相対パスで書きます。
\pt{読み方}2章で何をしたかとモデルの構成をつかみ、分からない言葉は3章の用語集で調べてください。4章で図の見方、5章で表の見方、6章で結論と、これからの選択肢を説明します。\textbf{結論だけを知りたいときは、6章から読んでください}。
\pt{元のノート}\file{01_機械学習コードフォルダ} の \file{材料別_温度シフトGPR_ARD対応Matern52＋線形＋ホワイトノイズ.ipynb} です。6月のノート（\file{GPR6月_ARD対応Matern52カーネル ＋ ホワイトノイズ.ipynb}）を基準にしています。実行は約 @@RUN_MIN@@ 分です。
\pt{記号表}記号・モデル名・列名の一覧は、\file{results/学習結果説明_記号表.pdf} の「材料別 GPR」の章にあります。
\pt{作り直し方}\file{05_材料ごとの機械学習データと予測検証データ} フォルダで \texttt{lualatex 学習結果説明\_材料別温度シフトGPR.tex} を2回実行すると、この PDF を作り直せます。
\pt{ノートの判定}ノートの各節の最後に、\file{>> 判定：OK}（想定どおり）か \file{>> 判定：★要確認}（想定と違う、または注意が要る点）が出ます。今回は、OK が @@N_OK@@ 個、★ が @@N_STAR@@ 個でした。★ はエラーではなく、結果を読むときに気をつける点の目印です。6章に、★ の意味を書きました。
\end{points}

% =====================================================================
\section{何をしたか（概要）}

\subsection{目的と、モデルの考え方}
射出成形の金型の中を流れる樹脂の見かけの粘度 $\eta$ を、せん断速度 $\gd$（流れの速さの目安）、圧力 \PIII、流動中樹脂温度 $\Tff$、加熱シリンダ温度 $\Tset$ から予測するモデルを、材料ごとに作りました。Cross-WLF 式（粘度の式としてよく使われる式）と同じく、「温度が変わると、粘度の曲線がずれる（シフトする）」という構造を、ガウス過程回帰の中に入れました。このずれの大きさを表すのが\textbf{シフトファクター} $a(T)$ です。

ずれ方は 2 通りあります。\textbf{横のシフト} $a_h$ は「温度が変わると、同じ粘度になるせん断速度が $a_h$ 倍ずれる」、\textbf{縦のシフト} $a_v$ は「同じせん断速度で、粘度が $a_v$ 倍になる」です。どちらも基準温度 $T_{\mathrm{ref}}$ で 1 です。本体の検討（手順5a）で、横と縦を別々に決める形（構造 B）が、データに最もよく合いました。今回もこの形を使いました。

\subsection{データ}
\begin{table}[H]
\caption{使ったデータ（問題データ一覧の「除外」を除いた後）。5data で学習し、3data で検証する}\label{tab:data}
\small
\begin{tabularx}{\linewidth}{@{}lrrrrX@{}}
\toprule
材料 & 学習の行 & 学習の条件 & 検証の行 & 検証の条件 & 設定温度・厚み \\
\midrule
@@TABLE_DATA@@
\bottomrule
\end{tabularx}
\end{table}
検証の 3data は、学習の 5data の射出速度の範囲（10〜240 mm/s）の\textbf{外側}です。低速側（3・5 mm/s。厚い板だけ）と高速側（280・320 mm/s）に分けて評価しました。つまり、3data の予測は「外挿」（学習した範囲の外への予測）です。

PP の 3data のファイルには、加熱シリンダ温度の列がありません。ノートは、元ファイル名（例：\file{071_190℃_280_0.5mm.csv}）から設定温度を読み取ります。

\subsection{モデルの構成}
\begin{table}[H]
\caption{比べた 3 つのモデル。$f$ はガウス過程。$\ln$ は自然対数。\PIII\ は Pa に直してから $\ln$、$\Tff$ は K に直してから $\ln$}\label{tab:model}
\small
\begin{tabularx}{\linewidth}{@{}l>{\raggedright\arraybackslash}X@{}}
\toprule
モデル & 式と考え方 \\
\midrule
B\_Tset & $\ln\eta=\ln a_v(\Tset)+f(\ln\gd+\ln a_h(\Tset),\ \ln\PIII,\ \ln\Tff)$。設定温度でずらす（本体の「モデル上の溶融温度＝設定温度」と同じ） \\
B\_Tff & $\ln\eta=\ln a_v(\Tff)+f(\ln\gd+\ln a_h(\Tff),\ \ln\PIII,\ \Tset)$。流動中樹脂温度でずらす（Cross-WLF と同じく、実際に流れている樹脂の温度） \\
C（基準） & $\ln\eta=f(\ln\gd,\ \ln\PIII,\ \ln\Tff,\ \Tset)$。シフトの構造なし。4 つを並べただけの通常の GPR（6月のノートの形に $\Tset$ を足したもの） \\
\midrule
カーネル & 6月と同じ ARD Matern 5/2 ＋ ホワイトノイズに、線形項（DotProduct）を足した（3data は外挿なので）。ノイズ分散の下限は $10^{-6}$ \\
シフトの式 & PP・粉砕材混合 PP：アレニウス型 $\ln a=(E/R)(1/T-1/T_{\mathrm{ref}})$、$T_{\mathrm{ref}}=210$ ℃。係数は $\Eh$・$\Ev$（kJ/mol）。GPPS：WLF 型 $\ln a=-A_1[(T-D_2)/(A_2+T-D_2)-(T_{\mathrm{ref}}-D_2)/(A_2+T_{\mathrm{ref}}-D_2)]$、$A_2=51.6$ K・$D_2=T_g$（仮置き 100 ℃）を固定、$T_{\mathrm{ref}}=220$ ℃。係数は $A_{1,h}$・$A_{1,v}$ \\
入力の取り方 & @@INPUT_MODE_TEXT@@ \\
\bottomrule
\end{tabularx}
\end{table}

\subsection{推定と評価の流れ}
\begin{enumerate}[itemsep=0.3ex,topsep=0.5ex]
\item \textbf{外れ値の検知}：これまでと同じ規則 f（条件平均が、ほかの条件からの予測と大きく違う条件）で調べ、PP・粉砕材混合 PP で拾った条件を、\file{05_材料ごとの機械学習データと予測検証データ/問題データ一覧.csv} に「除外」で登録して除く。
\item \textbf{シフトファクターの係数}（$\theta_h$・$\theta_v$）を、\textbf{プロファイル尤度}で決める：係数の格子の各点で GPR を学習し、データへの合い具合（対数周辺尤度 LML）が最大の点を採る。粗い格子 → 最良の点を中心にした細かい格子 →（最良の点が動いたら）取り直し、の順。
\item \textbf{念入りな最適化}：最良の係数と $\theta=0$（シフトなし）の 2 点だけを、最適化のやり直し 20 回で学習し直す。ハイパーパラメータの尤度の面に局所解が多く、格子の各点（やり直し 1 回）の LML は 1〜3 ほどずれることがあるため。$\theta=0$ との LML の差（$\Delta$LML）は、この 2 つから求める。最良の係数の学習が、そのまま最終モデルになる（基準の C も、同じ回数で学習する）。
\item 最終モデルで 3data を予測して、誤差と予測区間を評価する。参考に、6月のノートと同じ Cross-WLF 式とも比べる。
\item \textbf{補助の確認}（ノートの外の \file{04_補助スクリプト/mat_seed_check.py}）：乱数シード（最適化の初期値）を変えて学習し直し、LML と 3data の RMSE がどれだけばらつくかを見る。
\end{enumerate}

\subsection{外れ値と、問題データ一覧の更新}
@@OUTLIER_TEXT@@

% =====================================================================
\section{用語集}
\begin{termlist}
\term{ガウス過程回帰（GPR）}入力と出力の関係を、決まった式の形を仮定せずに、データから滑らかな関数として学ぶ方法。予測の平均と一緒に、予測のばらつき（不確かさ）も出せる。
\term{カーネル}「入力が近い点どうしは、出力も似ている」という度合いを決める関数。ARD Matern 5/2 は、入力ごとに「どれだけ離れると似なくなるか」（長さスケール）を別々に学ぶ。線形項は、学習の範囲の外を一次式で延ばす役目。ホワイトノイズは、同じ条件でもショットごとにばらつく分。
\term{長さスケール・特徴スケール}入力がどれだけ変わると、出力が大きく変わるか。小さいほど、その入力に敏感。6月のノートの「重要度」は、その逆数（1／長さスケール）。
\term{シフトファクター $a(T)$}温度による粘度の曲線のずれ。横（$a_h$）と縦（$a_v$）がある。$T_{\mathrm{ref}}$ で 1。
\term{アレニウス型・WLF 型}シフトファクターの温度の式。結晶性の樹脂（PP）はアレニウス型、非晶性の樹脂（GPPS）は WLF 型がよく使われる。
\term{見かけのシフトエネルギー $\Eh$・$\Ev$}アレニウス型の係数（kJ/mol）。文献の流動の活性化エネルギーとは別物で、この「見かけの粘度」の温度依存性を表す、データから決めた値。
\term{対数周辺尤度（LML）}モデルがデータにどれだけ合うかを表す数（大きいほど合う）。複雑すぎるモデルには自動で罰がかかる。
\term{プロファイル尤度}係数（ここではシフトの係数）を格子で振り、各点でほかのパラメータを最適にしたときの LML。最大から 1.92 以内の範囲が、おおよその 95\% 区間。
\term{尾根}プロファイル尤度の山が、斜めの帯になること。横の係数 $\theta_h$ と縦の係数 $\theta_v$ は別々には決まりにくく、尾根に沿って尤度がほとんど変わらない。差 $\theta_v-\theta_h$ が最もよく決まる。
\term{局所解・最適化のやり直し}ハイパーパラメータを決める最適化は、出発点によって、別々の「谷底」（局所解）に着くことがある。出発点を乱数で変えて何回かやり直し、LML が最大の解を採る。やり直しの回数が少ないと、LML の低い谷底で止まることがある。
\term{乱数シード}乱数の初期値の決め方。シードを変えると、最適化の出発点が変わる。同じシードなら、同じ結果になる。
\term{$\theta=0$}シフトなし。B\_Tset では「$\Tset$ を使わない」、B\_Tff では「$\Tff$ を使わない」モデルと同じ。
\term{$\Delta$LML・尤度比検定の 5\% 基準}最良の係数と $\theta=0$ の LML の差。2 つの係数を足した効果が認められる目安は、$2\Delta\mathrm{LML}>5.99$、つまり $\Delta\mathrm{LML}\ge3.0$。
\term{マスターカーブ}温度ごとの曲線を、シフトファクターでずらして重ねた曲線。重なれば、温度の効果がシフトで表せている。
\term{条件・条件平均}条件は、設定温度・厚み・射出速度の組。条件平均は、同じ条件のショット（4〜5 回）の値の平均。
\term{外挿・内挿}学習した範囲の外への予測（外挿）と、中への予測（内挿）。外挿は一般に難しい。
\term{RMSE}予測の誤差の大きさ（$\ln\eta$ の二乗平均平方根）。0.1 は粘度で約 10\%。
\term{決定係数 $R^2$}1 に近いほどよく当たる。$\ln\eta$ で計算したものと、6月のノートと同じく粘度の実数値で計算したもの（R²（実数））を出す。実数値の $R^2$ は、粘度の大きい点（低速側）の誤差に強く左右される。
\term{95\% 予測区間・被覆率}新しく測った値が 95\% の確率で入るはずの範囲と、実際に入った割合。0.95 に近いほど、ばらつきの見積もりが正しい。
\term{NLPD}予測分布の当たり具合（平均の誤差と区間の幅の両方を見る）。小さいほどよい。
\term{リーク}目的変数を計算する式に入っている量を説明変数に使うと、モデルが「式」を覚えてしまい、見かけだけ精度が高くなること。
\term{規則 f}外れ値の検知の規則。条件を 1 つずつ抜いて、残りの条件から予測し、予測と実測の差が大きい（$|z|>3.5$ かつ差が約 10\% 超）条件を拾う。
\term{問題データ一覧}問題のあるデータを 1 行ずつ書いた表（\file{05_材料ごとの機械学習データと予測検証データ/問題データ一覧.csv}）。区分「除外」の行は、ノートが除く。
\end{termlist}

% =====================================================================
\section{図の見方}

図は材料ごとのフォルダにあります。ここでは主に PP の図で見方を説明し、ほかの材料の結果は文中と5章の表で示します。

@@FIGURES@@

% =====================================================================
\section{表の見方}

@@TABLES@@

% =====================================================================
\section{結論と、これからの選択肢}

@@CONCLUSION@@

\end{document}
